\documentclass[11pt,a4paper]{article}
\usepackage[T1]{fontenc}
\usepackage[utf8]{inputenc}
\usepackage{amsmath,amssymb,amsthm}
\usepackage[margin=1in]{geometry}
\usepackage[colorlinks=true,linkcolor=blue,urlcolor=blue]{hyperref}
\theoremstyle{plain}
\newtheorem{theorem}{Theorem}
\newtheorem{lemma}[theorem]{Lemma}
\newtheorem{proposition}[theorem]{Proposition}
\newtheorem{corollary}[theorem]{Corollary}
\theoremstyle{definition}
\newtheorem{remark}[theorem]{Remark}

\title{The empirical recurrence for OEIS A211556, proved by a one-dimensional reduction}
\author{Adrian Perez Fontelles\\ \small Independent researcher}
\date{4 September 2026}
\begin{document}
\maketitle

\begin{abstract}
OEIS A211556 counts symmetric $(n+1)\times(n+1)$ matrices over
$\{-9,\dots,9\}$ in which every $2\times2$ subblock sums to zero and has two, three or four distinct
entries. Both sides of the matrix grow with $n$, so there is no slice of bounded width to walk
along. There does not need to be: the sum-zero condition is linear and solves completely, and
its solutions are indexed by a SEQUENCE --- the diagonal of the matrix. The count becomes a walk
on $S=3160$ states, and the empirical recurrence of order $20$ contributed by R.~H.~Hardin
follows.
\end{abstract}

\noindent\small 2020 Mathematics Subject Classification. 05A15, 15B36, 15A18.\normalsize

\section{The conjecture}

OEIS A211556 is ``Number of (n+1)X(n+1) -9..9 symmetric matrices with every 2X2 subblock having sum zero and two, three or four distinct values.''. It has offset $1$ and begins
\[
180, 1712, 16318, 155904, 1492928, 14327826, 137797604, 1327935624,\ \dots
\]
The entry states, as a conjecture contributed by its author and never marked settled:
\begin{quote}
Empirical: a(n) = 94*a(n-1) -3990*a(n-2) +100725*a(n-3) -1674425*a(n-4) +19135732*a(n-5) -151829733*a(n-6) +819081545*a(n-7) -2798053008*a(n-8) +4813088305*a(n-9) +909604521*a(n-10) -14434966072*a(n-11) -404287970*a(n-12) +26746796493*a(n-13) +30215600922*a(n-14) +16403931314*a(n-15) +5216836244*a(n-16) +1016849976*a(n-17) +119142480*a(n-18) +7668288*a(n-19) +207360*a(n-20)
\end{quote}
As of the ``Last modified'' line on the live entry (Oct 03 2025 16:13:43, revision 8) this is still
recorded as empirical, and nothing on the entry records it as proved.

\section{The sum-zero condition solved}

Let $x$ be an $N\times N$ symmetric matrix with entries in $\{-9,\dots,9\}$ such that
\[
x(i,j)+x(i,j{+}1)+x(i{+}1,j)+x(i{+}1,j{+}1)=0
\]
for every $0\le i,j\le N-2$.

\begin{lemma}\label{lem:solve}
Such matrices are exactly the matrices
\[
x(i,j)=(-1)^{i+j}\,\frac{u_i+u_j}{2},
\qquad u_0,\dots,u_{N-1}\in\{-9,\dots,9\}\ \text{all of the same parity},
\]
and the sequence is recovered from the matrix by $u_i=x(i,i)$, so the correspondence is a
bijection.
\end{lemma}

\begin{proof}
Put $y(i,j)=(-1)^{i+j}x(i,j)$. The four entries of a $2\times2$ subblock alternate in sign
under this substitution, so
\[
x(i,j)+x(i,j{+}1)+x(i{+}1,j)+x(i{+}1,j{+}1)
=(-1)^{i+j}\bigl[y(i,j)-y(i,j{+}1)-y(i{+}1,j)+y(i{+}1,j{+}1)\bigr],
\]
and the condition says exactly that the mixed second difference of $y$ vanishes everywhere.
Summing that identity over a rectangle gives
$y(i,j)=y(i,0)+y(0,j)-y(0,0)$ for all $i,j$, which is $y(i,j)=f(i)+g(j)$ with
$f(i)=y(i,0)-y(0,0)$ and $g(j)=y(0,j)$; conversely every such $y$ has vanishing mixed
difference. Symmetry of $x$ gives $y(i,j)=y(j,i)$, hence $f(i)-g(i)=f(j)-g(j)$ for all $i,j$:
the difference is a constant $c$. Writing $u_i=2\bigl(f(i)-c/2\bigr)$, which is an integer
because $y(i,i)=x(i,i)$ is, we get $y(i,j)=(u_i+u_j)/2$ as claimed. Taking $j=i$ gives
$u_i=y(i,i)=x(i,i)$.

For the two side conditions: $u_i+u_j$ must be even for all $i,j$, which says all the $u_i$ have
the same parity; and $|x(i,j)|=|u_i+u_j|/2\le9$ for all $i,j$ is equivalent to $|u_i|\le9$
for each $i$, the case $j=i$ giving one direction and the triangle inequality the other.
\end{proof}

So the matrices of size $N$ correspond one for one to the sequences $u_0,\dots,u_{N-1}$ in
$\{-9,\dots,9\}$ with all entries of one parity. What remains is to say, in terms of $u$,
that every $2\times2$ subblock has two, three or four distinct entries.

\begin{lemma}\label{lem:block}
Write $\pi_i=(u_i,u_{i+1})$ for the $i$-th consecutive pair. The number of distinct entries of
the subblock at $(i,j)$ depends only on $\pi_i$ and $\pi_j$: with $\pi_i=(\alpha,\beta)$ and
$\pi_j=(\gamma,\delta)$ it is
\[
\Bigl|\Bigl\{\tfrac{\alpha+\gamma}{2},\,-\tfrac{\alpha+\delta}{2},\,
-\tfrac{\beta+\gamma}{2},\,\tfrac{\beta+\delta}{2}\Bigr\}\Bigr| .
\]
\end{lemma}

\begin{proof}
By Lemma~\ref{lem:solve} the four entries are $(-1)^{i+j}$ times those four numbers, in
order. Multiplying all four by $\pm1$ does not change how many are distinct.
\end{proof}

Call two consecutive pairs COMPATIBLE when that number lies in $\{2,3,4\}$. The condition on the
matrix is that $\pi_i$ and $\pi_j$ are compatible for every $i$ and $j$ --- including $i=j$,
which asks each pair to be compatible with itself. In other words:

\begin{corollary}
The matrices counted are exactly the sequences $u$ whose set of consecutive pairs is a clique in
the compatibility graph on the $\le 361$ pairs.
\end{corollary}

\section{A walk on the sequence}

The clique condition is not local: a pair must be compatible with every pair anywhere else in
the sequence. Carrying the set of pairs already used would be an exponential state. Carrying
what it FORBIDS is not.

\begin{lemma}\label{lem:lump}
Let $A$ be the set of pairs compatible with every pair used so far. Then a sequence may be
continued exactly by pairs lying in $A$, and appending a pair $\pi$ replaces $A$ by
$A\cap N(\pi)$, where $N(\pi)$ is the set of pairs compatible with $\pi$. Hence the pair
$(\,\text{last value of }u,\;A\,)$ carries everything the future needs.
\end{lemma}

\begin{proof}
The condition is a conjunction of pairwise compatibilities, so the set of pairs still allowed
after using a set $T$ is $\bigcap_{\pi\in T}N(\pi)$, which is a function of $T$ and is updated
by intersection as $T$ grows. The next pair is $(\text{last value},w)$, so the last value is
the only other thing a step needs.
\end{proof}

The sets $A$ that arise are intersections of neighbourhoods, of which there are far fewer than
subsets; enumerated from the start $A_0=\{\pi:\pi\in N(\pi)\}$, the reachable states number
$S=3160$. The two parity classes never mix, so the walk starts by choosing a parity and a first
value, and every state may end a sequence:
\[
a(n)\;=\;\iota^{\!\top}M^{\,j}\tau ,\qquad \tau=\mathbf{1},\qquad S=3160.
\]
The walk index $j$ is the size of the matrix, and it differs from the entry's own $n$ by a
constant read off by matching the published terms rather than assumed. In particular $a$ is
$C$-finite --- for a family in which both dimensions grow with $n$.

\section{The proof}

\begin{lemma}\label{lem:crit}
Let $q(t)=t^{20}-\sum_i c_it^{20-i}$ be the characteristic polynomial of the
conjectured recurrence. The residual $a(n)-\sum_ic_ia(n-i)$ equals
$\iota^{\!\top}M^{\,j}q(M)\tau$ for the corresponding walk index $j$, so the recurrence
holds from some point on if and only if $\iota^{\!\top}M^{j}q(M)\tau=0$ for all large $j$,
and it is enough to examine $j<S$.
\end{lemma}

\begin{proof}
Factor $M^{j}$ out of $M^{j+20}-\sum_ic_iM^{j+20-i}$. For the bound, $M^{S}$
is an integer combination of $I,M,\dots,M^{S-1}$ by Cayley--Hamilton, so the numbers
$u_j=\iota^{\!\top}M^{j}q(M)\tau$ obey the monic recurrence given by the characteristic
polynomial of $M$; $S$ consecutive zeros force every later one, and the last nonzero $u_j$
pins the threshold exactly.
\end{proof}

\begin{theorem}
\[
a(n)\;=\;94\,a(n-1) - 3990\,a(n-2) + 100725\,a(n-3) - 1674425\,a(n-4) + 19135732\,a(n-5) - 151829733\,a(n-6) + 819081545\,a(n-7) - 2798053008\,a(n-8) + 4813088305\,a(n-9) + 909604521\,a(n-10) - 14434966072\,a(n-11) - 404287970\,a(n-12) + 26746796493\,a(n-13) + 30215600922\,a(n-14) + 16403931314\,a(n-15) + 5216836244\,a(n-16) + 1016849976\,a(n-17) + 119142480\,a(n-18) + 7668288\,a(n-19) + 207360\,a(n-20)
\]
for every $n>19$.
\end{theorem}

\begin{proof}
The vector $w=q(M)\tau$ was formed by $20$ matrix--vector products in exact integer
arithmetic and $\iota^{\!\top}M^{j}w$ evaluated for $j=0,1,\dots$, again exactly; those
integers vanish from the index corresponding to $n=19+1$ onwards and the run continued
until the working vector was identically zero, which settles every larger $j$ at once.
\end{proof}

\section{Verification}

Four checks, each independent of the algebra above.

First, the model reproduces all $16$ terms the entry publishes, exactly, in integer
arithmetic.

Second, the count was recomputed for the smallest matrices straight from the definition:
symmetric matrices were written out entry by entry over $\{-9,\dots,9\}$ and each
$2\times2$ subblock tested for summing to zero and for having two, three or four distinct entries. The
sequence $u$, the sign pattern, the clique and the state play no part in that computation, and
the published terms came back. Lemma~\ref{lem:solve} was separately checked the same way: for
small sizes and small ranges the number of symmetric sum-zero matrices found by direct
enumeration equals the number of admissible sequences.

Third, the conjectured recurrence was evaluated directly on the published terms, with no
matrices involved, and holds wherever the proved range and the data overlap.

Fourth, the annihilation test of Lemma~\ref{lem:crit} was carried out in exact integer
arithmetic on the whole vector rather than on a sampled prefix.

\begin{thebibliography}{9}
\bibitem{oeis} The OEIS Foundation, \emph{The On-Line Encyclopedia of Integer
Sequences}, \url{https://oeis.org}, 2026. Sequence A211556.
\bibitem{stanley} R.~P.~Stanley, \emph{Enumerative Combinatorics}, Volume 1, 2nd ed.,
Cambridge University Press, 2011. (Transfer-matrix method, Section 4.7.)
\bibitem{kemeny} J.~G.~Kemeny and J.~L.~Snell, \emph{Finite Markov Chains}, Springer, 1976.
(Lumping, Section 6.3.)
\bibitem{horn} R.~A.~Horn and C.~R.~Johnson, \emph{Matrix Analysis}, 2nd ed., Cambridge
University Press, 2013.
\end{thebibliography}

\end{document}
